\section{Methodology}
\label{sec:method}

\subsection{Mask classification and the EoMT model}
All segmentation models in this work follow the mask-classification
paradigm~\cite{maskformer}: the network predicts a fixed set of $N$ object
queries, each producing (i) a class distribution over $C{+}1$ labels (the
$C$ task classes plus a ``no-object'' class) and (ii) a binary mask over the
image. EoMT~\cite{eomt} realises this with a \emph{plain} DINOv2~\cite{dinov2}
Vision Transformer (we use \texttt{vit\_base\_patch14\_reg4\_dinov2}). A set
of $N$ learnable query embeddings is concatenated to the patch tokens and
processed jointly inside the last $L$ encoder blocks (here $L{=}3$), so the
same self-attention layers refine both image features and queries; no
separate pixel or Transformer decoder is used. A linear \emph{class head}
maps each query to its class logits, while a three-layer MLP \emph{mask head}
produces a query embedding that is multiplied with up-scaled patch features
to yield the per-query mask. During training EoMT uses \emph{mask-attention
annealing}: early on, query-to-patch attention is constrained to the
currently predicted mask (as in Mask2Former~\cite{mask2former}), and a
per-block probability schedule gradually relaxes this constraint so that, at
convergence, the plain encoder attention is used at inference.

\paragraph{Semantic inference.}
For a pixel $(h,w)$ a class score is obtained by marginalising the queries,
following MaskFormer:
\begin{equation}
  z_c(h,w) \;=\; \sum_{i=1}^{N} p_i(c)\, \sigma\big(m_i(h,w)\big),
  \label{eq:seminf}
\end{equation}
where $p_i(c)$ is the (no-object-excluded) class probability of query $i$ and
$\sigma(m_i)$ its sigmoid mask activation; the label is $\arg\max_c z_c$. This
same rule turns the panoptic COCO model into a semantic predictor, which is
essential for the comparison in \cref{sec:res-task4}.

\subsection{Training objective and optimisation}
Queries are matched to ground-truth segments by Hungarian bipartite
matching, and the matched pairs are supervised by a weighted sum of a
classification cross-entropy and a mask loss combining binary cross-entropy
and Dice terms. We use the EoMT defaults: class, mask and Dice coefficients
$2{:}5{:}5$ and a no-object weight of $0.1$. Optimisation uses AdamW with
weight decay $0.05$ and \emph{layer-wise learning-rate decay} (LLRD, factor
$0.8$), so deeper backbone blocks receive exponentially smaller learning
rates---this protects the pre-trained DINOv2 features while still allowing the
task-specific heads to move quickly. All training uses automatic mixed
precision (AMP).

\subsection{Task 4: a shared protocol for two label spaces}
\label{sec:method-task4}
We are given two EoMT checkpoints: one trained on Cityscapes (19 semantic
classes) and one on COCO (133 panoptic classes). Comparing them fairly is
non-trivial because the label spaces and tasks differ. We evaluate \emph{both}
as semantic segmenters on the full Cityscapes validation set using one
identical pipeline (\cref{eq:seminf}, argmax, per-class IoU). The Cityscapes
model predicts the 19 train-ids directly. The COCO model's logits are
projected into Cityscapes space through a fixed COCO$\rightarrow$Cityscapes
mapping (e.g.\ COCO \texttt{road}/\texttt{pavement} $\rightarrow$ Cityscapes
\texttt{road}/\texttt{sidewalk}, \texttt{tree}/\texttt{grass}
$\rightarrow$ \texttt{vegetation}). When several COCO classes map to the same
Cityscapes class we take the \emph{maximum} logit rather than the sum: this
avoids unfairly inflating Cityscapes classes that happen to be covered by many
COCO categories, giving a class high confidence whenever \emph{any} of its
COCO counterparts is strongly predicted. Three Cityscapes classes without a
reliable COCO correspondence---\texttt{pole}, \texttt{traffic sign} and
\texttt{rider}---are excluded from both models, so the reported mIoU is
computed over the remaining \textbf{16 classes}. COCO categories with no
Cityscapes equivalent (e.g.\ \texttt{pizza}) are routed to the ignore label.
The two networks use different native pre-processing (Cityscapes:
$1024$-px sliding windows; COCO: $640$-px resize-and-pad), but the metric,
mapping and exclusion list are shared.

\subsection{Task 5: staged COCO\,$\rightarrow$\,Cityscapes fine-tuning}
\label{sec:method-task5}
The COCO model carries broad visual priors but has never seen the Cityscapes
label space. Following the project guidance to ``first fine-tune the
prediction head, then gradually unfreeze'', we adopt a three-stage schedule
that monotonically increases the trainable capacity, all at $512{\times}512$
resolution with $N{=}200$ queries:
\begin{itemize}\setlength\itemsep{1pt}
  \item \textbf{Stage~0} (10 epochs, lr $10^{-4}$): train only the
  \texttt{class\_head} and \texttt{mask\_head}; DINOv2, queries and the
  up-scaling blocks stay frozen. The class head is re-initialised because
  Cityscapes has a different number of classes than COCO.
  \item \textbf{Stage~1} (15 epochs, lr $10^{-4}$): keep DINOv2 frozen but
  unfreeze the \emph{full} prediction head (queries, class head, mask head and
  up-scaling), initialised from Stage~0.
  \item \textbf{Stage~2} (20 epochs, lr $10^{-5}$, LLRD $0.8$): additionally
  unfreeze the last two DINOv2 blocks, initialised from Stage~1 with a fresh
  optimiser so the polynomial schedule restarts cleanly. The low learning rate
  and LLRD prevent the backbone features from being destroyed.
\end{itemize}
The fine-tuned checkpoints are evaluated with exactly the Task~4 protocol
(16-class mIoU) so the new model, the provided Cityscapes model and the COCO
model are directly comparable. We deliberately avoid a backbone-token cache
for the head-only stage: caching tokens from one augmented view while masks
are reloaded from another misaligns supervision and collapses the mask/Dice
loss.

\subsection{Post-hoc anomaly scoring}
\label{sec:method-anomaly}
Anomaly segmentation is performed without retraining, by converting a
segmentation network's outputs into a per-pixel anomaly score. From the
semantic logits $z_c(h,w)$ (\cref{eq:seminf}) with temperature $T$ we use:
\begin{align}
  s_{\text{MSP}} &= 1-\max_c \mathrm{softmax}(z/T)_c, \\
  s_{\text{MaxLogit}} &= -\max_c z_c/T, \\
  s_{\text{Entropy}} &= -\textstyle\sum_c \mathrm{softmax}(z/T)_c \log \mathrm{softmax}(z/T)_c .
\end{align}
For the mask architecture we additionally use Rejected-by-All
(RbA)~\cite{rba}: a pixel is anomalous when \emph{no} query both predicts a
known class confidently and covers that pixel, implemented as
$s_{\text{RbA}} = 1-\max_i\big[\sigma(m_i(h,w))\cdot\max_{c}p_i(c)\big]$.
Temperature scaling, a simple calibration technique, is applied to the
logit-based scores to test whether sharpening or smoothing the distribution
improves separability. Anomaly quality is measured by the Area under the
Precision--Recall curve (AuPRC, anomalies as positives) and the false-positive
rate at $95\%$ true-positive rate (FPR$_{95}$); higher AuPRC and lower
FPR$_{95}$ are better.
